El primer paso fue conseguir el código exacto de los autores. Lo publicaron en Zenodo, un repositorio científico, con un identificador permanente.

Lo bajamos y comprobamos su MD5, una huella del archivo. Si coincide, sabemos que es exactamente lo que publi- caron, sin un byte cambiado.

¿Por qué no programarlo nosotros? Porque entonces mediríamos nuestra forma de programarlo. Usando su artefacto, medimos su código, el de Folly y el de Anderson tal como ellos los usaron.

Un detalle práctico: el artefacto pide la versión del día de dos bibliotecas de pruebas. Esa versión ya cambió, así que fijamos dos versiones concretas para que se pueda repetir.
